% Post-CP21 work. All geometric hypotheses below are explicit.
\subsection{A rank-drop-compatible three-theta graph reduction}
\status{Candidate; the graph and closure arguments are complete, while the stated geometric input remains conditional.}

\begin{lemma}[Clopen closure of an invariant Lang subgroup]
\label{lem:normal-base-clopen-Lang-closure-candidate}
Let $B$ be normal, integral and Noetherian, let $E$ be a vector bundle
with Frobenius-semilinear map, and let $\mathcal L=\{Fv=v\}$ be its
Lang scheme. Suppose $U\subset B$ is a nonempty open where the stable
rank is constant, and let $\mathcal S_U\subset\mathcal L_U$ be a
finite \etale subgroup local system arising from an invariant
$\mathbf F_2$-subspace. Its closure $\overline{\mathcal S}$ in
$\mathcal L$ is open and closed and is a subgroup scheme.
If $C$ is integral and connected, $C\to B$ meets $U$ generically,
and $T\hookrightarrow\mathcal L_C$ is a finite \etale subgroup,
then $T\cap\overline{\mathcal S}_C$ is finite \etale of constant
rank. If $\mathcal S$ is generically contained in $T$, its entire
closure on $C$ is contained in $T$ and retains its generic rank.
If $T$ is generically contained in $\mathcal S$, it is contained
in the closure on every fibre.
\end{lemma}
\begin{proof}
The Lang equations have derivative $-I$, so $\mathcal L\to B$ is
separated, quasi-finite and \etale. It is normal and Noetherian,
and its finitely many irreducible components are disjoint and open
and closed. Each component dominates $B$: an \etale map is open,
and every nonempty open of the integral base contains its generic
point. Its intersection with $U$ is therefore a nonempty irreducible
open. Since $\mathcal S_U$ is open and closed in $\mathcal L_U$,
its closure is a union of entire components of $\mathcal L$.
The subgroup property follows by checking addition and inverse on
the dense inverse image of $U$; every component of
$\mathcal L\times_B\mathcal L$ also dominates $B$.

A finite \etale monomorphism $T\to\mathcal L_C$ is an open immersion,
and is also closed because it is finite and the target is separated
over $C$. Its intersection with the clopen closure is consequently
finite \etale, with locally constant rank on connected $C$.
If the closure has a point outside $T$, that open complement is
\etale over $C$, so its image meets the generic point, contradicting
the assumed generic containment. The reverse containment follows
from the same argument on the complement inside finite \etale $T$.
These arguments do not identify an entire rank-dropping Lang fibre
with the generic fibre.
\end{proof}

\begin{lemma}[The three-component graph with odd labels]
\label{lem:three-theta-odd-labelled-graph-candidate}
Let $K$ be a finite abelian group of odd order $k>1$, and put
$H=(K\setminus\{0\})/\{h\sim-h\}$. Form a graph with three vertices,
two loops $\ell_{i,h,\epsilon}$ for each vertex $i$ and type $h\in H$,
and two edges $e_{ij,g,\epsilon}$ for each $i<j$, $g\in K$,
$\epsilon\in\{0,1\}$. Let $G$ be its cochain space modulo vertex
coboundaries, so $\dim G=9k-5$.
For each $r=(r_1,r_2,r_3)$ with $r_1+r_2+r_3=0$, retain only the
six edges with $g=r_j-r_i$, and let $T_r\subset G$ be the kernel
of restriction to this doubled triangle. Then $\dim T_r=9k-9$,
and it has a direct sum decomposition into eighteen-dimensional
coordinate spaces $T_{r,h}$, consisting of the six loops of type
$h$ and the twelve cross edges with $g=r_j-r_i\pm h$.
Write $T_{r,h}^0$ for the coordinate-sum-zero hyperplane.
Then:
\begin{enumerate}
\item The span of all $T_{r,h}^0$, for all permitted $r,h$, is the
codimension-one subspace $G^0\subset G$ on which the sum of all edge
coordinates is zero.
\item If $S\subset G$ lies in every $T_r$ and is invariant under
$\prod_h S_{18}$ acting on the coordinate spaces of $T_0$, then $S=0$.
\end{enumerate}
\end{lemma}
\begin{proof}
There are $9k-3$ edges and two independent vertex coboundaries.
The selected doubled triangle has first cohomology dimension four.
The unselected edge cochains are independent modulo coboundaries:
a coboundary vanishing on its connected selected triangle is zero.
They therefore give precisely the stated coordinates of $T_r$.

Pair cochains modulo coboundaries with cycles. A cycle annihilating
$T_{r,h}^0$ has the same coefficient $c_{r,h}$ on all its eighteen
edges. The six loop edges of type $h$ are independent of $r$, so
$c_{r,h}=c_h$. Each individual difference $r_j-r_i$ ranges over all
of $K$ even with $\sum r_i=0$: assign these two entries $0,g$ and
the third $-g$. Thus every cross-edge coefficient equals $c_h$ for
every type $h$, forcing a single common coefficient $c$ on every
edge of the graph. The all-ones edge vector is a cycle, since the
cross degree at every vertex is $4k$, which is even. Hence this
annihilator is exactly one-dimensional, proving the first assertion.

For the second assertion, regard $S$ inside $T_0$. If any vector has
nonconstant coordinates in one type, subtract its image under a
transposition in that factor. This produces a pair vector in $S$;
conjugating by the whole factor gives every pair vector and hence
$T_{0,h}^0\subset S$. Choose such a pair vector supported on one
cross edge of type $h$ and one loop of the same type. Choose a
permitted $r$ which selects that cross edge. Its restriction to the
doubled triangle is nonzero, since precisely one of the two parallel
edges has value one. Thus this vector is not in $T_r$, a contradiction.
Consequently all vectors of $S$ are type-constant. Write one as
$\sum_h a_h\mathbf1_{T_{0,h}}$, and extend the coefficients to an
even function $a:K\to\mathbf F_2$ with $a(0)=0$.
For $r=(-u,0,u)$, restriction to the selected triangle vanishes only if
\[
                         a(u)+a(u)+a(2u)=0.
\]
Hence $a(2u)=0$ for every $u$. Multiplication by two is an
isomorphism of $K$, so all coefficients are zero.
\end{proof}

\begin{theorem}[Conditional odd-degree irreducibility and full image]
\label{thm:three-theta-conditional-full-image-candidate}
Let $A$ be a nonordinary principally polarized surface in
characteristic two and let $d>1$ be odd. Choose the polarized
isogeny $\phi:B\to A$ of degree $d^2$ used in
Lemma~\ref{lem:nonordinary-ppav-odd-sharp-norm-candidate}, and put
$K=\ker\phi$. Assume the following two additional geometric inputs.
\begin{enumerate}
\item A principal theta curve on $B$ has integral nodal norm, with
exactly two distinct transverse self-identifications for every
unordered nonzero kernel type and no triple identifications.
Three general theta translates with sum of translating points zero
have norms meeting transversely, with all their mutual intersections
and self-nodes distinct. Their independently $K$-translated source
unions with $r_1+r_2+r_3=0$ all lie in a connected coefficient open
where the additional norm identifications remain distinct ordinary
nodes, even though the source union itself is nodal.
\item On the smooth sharp norm family in the single source system
$|3\Theta_B|$, the node permutation image is the full product
$\prod_hS_{18}$, and, after a vector-space splitting, the boundary
image contains the actual $\GL_8$ block and every translation in
$\Hom(U,T_h^0)$ for each type $h$, where $\dim U=8$.
\end{enumerate}
Then the target odd-degree geometric and arithmetic images equal
$\GL_{9d^2-1}(\mathbf F_2)$.
\end{theorem}
\begin{proof}
Work geometrically and write $R=9d^2-1$, $\delta=9d^2-9$,
$t=(d^2-1)/2$. The pure norm sequence is
$0\to T\to V\to U\to0$, with $\dim U=8$.
Let $S$ be a nonzero proper global invariant subspace. Its projection
to $U$ on this norm family is either zero or all of $U$.
This choice is uniform on the connected dense maximal-rank part of
the source coefficient open:
the rank of a morphism of finite \etale local systems is locally
constant. The clopen closure argument below carries the resulting
alternative to the three-theta specializations; no constancy of
the rank-eight quotient is asserted at those rank-dropping points.
Their translating points have sum zero, so their source divisors
belong to this same linear system.
No inequality between $8$ and the number of node types is used.

Let $Y$ be the common target union of three normed theta translates.
The hypotheses identify its graph with the preceding lemma, and
identify the partial-normalization kernels with its $T_r$.
The graph cohomology embeds in the relative Lang fibre: ambient
$H^1(B,\mathbf F_2)$ restricts injectively to each principal theta
normalization, and therefore has zero intersection with the graph
part of the normalization sequence.
The relative coherent Frobenius bundle on the target projective
coefficient base is locally free of rank $R$, including $Y$:
for every ample divisor the ambient coherent $H^1$ injects and the
relative quotient has constant dimension, by the ample-vanishing
structure-sheaf sequence. Apply the clopen Lang closure lemma to
$S$ on the maximal-rank open and denote its special fibre by
$\overline S_Y$. It is the same subspace of the relative Lang fibre
at $Y$ for every source specialization. The graph subgroups
$T_r$ and their hearts are finite \etale on the source neighbourhoods
by the distinct-node hypothesis.

If the generic projection of $S$ to $U$ is zero, the closure lemma
puts $\overline S_Y$ inside every $T_r$, preserves its full dimension,
and makes it invariant under the product node group on $T_0$.
The second graph assertion gives $\overline S_Y=0$, a contradiction.

If the projection is surjective, the stated translations put every
$T_h^0$ into $S$ on the generic norm family. The closure lemma then
puts every $T_{r,h}^0$ into $\overline S_Y$, so its first graph
assertion gives $G^0\subset\overline S_Y$.
For a fixed $r$, the intersection $T_r\cap G^0$ has dimension
$\delta-1$: the total-edge-sum functional is nonzero on each
unselected-edge basis vector of $T_r$.
The finite \etale subgroup $T_r\cap\overline S$ has constant rank
on its connected norm-family base. Therefore
\[
            \dim(S\cap T)\ge\delta-1,
        \qquad \dim S\ge(\delta-1)+8=R-1.
\]
A proper $S$ must be a hyperplane. Over $\mathbf F_2$ its quotient
is a trivial one-dimensional representation, contradicting the
proved absence of invariant covectors in every odd nonordinary
principal degree. Thus the global representation is irreducible.

Let $J$ be its normal subgroup generated by transvections.
The nonzero $J$-socle is invariant under the full irreducible group,
so the representation restricted to $J$ is semisimple.
The actual natural $\GL_8$ block lies in $J$, and its natural
constituent occurs with multiplicity one. It therefore belongs to
one simple $J$-summand. All rank-one gluing maps into the hearts
place those hearts in the same summand. Its dimension is at least
\[
                       8+17t>\frac{8+18t}{2}=R/2.
\]
Its conjugates cannot be distinct disjoint summands, so it is
invariant under the full group and equals the whole representation.
Hence $J$ is irreducible. The characteristic-two transvection
classification, as applied in
Theorem~\ref{thm:mixed-every-fixed-modulus-all-odd-candidate}, leaves
only the full general linear group or a group preserving a
nondegenerate alternating form. The actual natural $\GL_8$ block
fixing a complement has no invariant bilinear form or functional,
so would lie in the radical of such a form. This excludes the latter
case, proving the full geometric image; the arithmetic image follows.
\end{proof}

\begin{remark}[The exact geometric gates retained]
\label{rem:three-theta-conditional-geometric-gates-candidate}
Neither additional hypothesis of the theorem is asserted uniformly
here. In particular, theta translates do not move their intrinsic
self-intersections with a fixed kernel translate. For higher odd
kernel order these intersections can be tangent, and theta norms
can also have triple identifications. Pairwise transversality alone
does not exclude the latter. Likewise an actual $\GL_8$ block does
not by itself supply the full product node groups or all gluing
hearts. The unconditional result currently available is the actual
block in every odd degree. This reduction isolates the remaining
geometric inputs while treating the common boundary's stable-rank
drop rigorously.
\end{remark}
